\cleardoublepage
\phantomsection
\begin{planbox}[แผนบริหารการสอนประจำบทที่ 4\quad งาน พลังงาน และหลักการเครื่องปั่นเหวี่ยง (สัปดาห์ที่ 4)]
\noindent\textbf{1. หัวข้อเนื้อหาประจำบท (บรรยาย 2 ชั่วโมง และปฏิบัติการ 3 ชั่วโมง)}
\begin{itemize}
    \item 4.1 งาน พลังงานจลน์ พลังงานศักย์ และกฎการอนุรักษ์พลังงานกล
    \item 4.2 กำลังทางฟิสิกส์และการคำนวณการใช้พลังงานของเครื่องมือในห้องปฏิบัติการจุลชีววิทยา
    \item 4.3 ฟิสิกส์ของเครื่องปั่นเหวี่ยง (Centrifuge) ความแตกต่างระหว่างความเร็วรอบ (RPM) กับแรงเหวี่ยงสัมพัทธ์ (RCF หรือ $\times g$)
    \item 4.4 ปฏิบัติการที่ 4 การอนุรักษ์พลังงานกลและการคำนวณตั้งค่าเครื่องปั่นเหวี่ยงพร้อมกฎการจัดสมดุลโรเตอร์
\end{itemize}

\noindent\textbf{2. วัตถุประสงค์เชิงพฤติกรรมและผลลัพธ์การเรียนรู้รายวิชา (CLOs ที่เชื่อมโยง PLOs)}
\begin{enumerate}
    \item คำนวณงาน พลังงาน และกำลังของอุปกรณ์ไฟฟ้าในห้องปฏิบัติการได้ถูกต้อง \textbf{(CLO1 เชื่อมโยง PLO2)}
    \item อธิบายเหตุผลทางฟิสิกส์ว่าทำไมการรายงานวิธีปั่นเหวี่ยงแยกเซลล์ในงานวิจัยจุลชีววิทยาต้องระบุค่า RCF ($\times g$) แทนค่า RPM และคำนวณแปลงค่าระหว่าง RPM กับ RCF ได้ \textbf{(CLO1, CLO2 เชื่อมโยง PLO2, PLO3)}
    \item ปฏิบัติการจัดวางหลอดตัวอย่างในหัวปั่นโรเตอร์ให้สมดุลตรงข้ามกันอย่างปลอดภัยตามหลักโมเมนต์ความเฉื่อย \textbf{(CLO3, CLO4 เชื่อมโยง PLO1, PLO5, PLO6)}
\end{enumerate}

\noindent\textbf{3. กิจกรรมการจัดการเรียนรู้เชิงรุก (Active Learning \& OBE)}
\begin{itemize}
    \item \textbf{Case-Based Learning}\quad ยกกรณีศึกษานักศึกษาคัดลอกวิธีทดลองจากบทความวิจัยที่ระบุ ``ปั่นเหวี่ยงที่ $4{,}000\text{ rpm}$'' มาใช้กับเครื่องปั่นขนาดเล็ก (Microcentrifuge) แล้วปรากฏว่าเซลล์แบคทีเรียไม่ยอมตกตะกอนลงก้นหลอด เพราะรัศมีโรเตอร์ต่างกัน
    \item \textbf{Hands-on Rotor Balancing}\quad ฝึกจัดเรียงหลอดไมโครทิวบ์จำนวน 2, 3, 4, 6 หลอด ในโรเตอร์ 12 ช่อง และ 24 ช่อง ให้แรงหนีศูนย์กลางหักล้างกันเป็นศูนย์
\end{itemize}

\noindent\textbf{4. สื่อการเรียนรู้และแหล่งเรียนรู้}
\begin{itemize}
    \item เอกสารคำสอนบทที่ 4 แผนภาพสมดุลโรเตอร์ และเครื่องปั่นเหวี่ยง Microcentrifuge ในห้องปฏิบัติการ
\end{itemize}

\noindent\textbf{5. การวัดและประเมินผลประจำบท}
\begin{itemize}
    \item การตรวจใบงานแปลงค่า RPM เป็น RCF \textbf{(CLO1, CLO2)}
    \item การประเมินทักษะการจัดสมดุลหลอดปั่นเหวี่ยงและรายงานปฏิบัติการที่ 4 \textbf{(CLO3, CLO4)}
\end{itemize}
\end{planbox}
\clearpage

\chapter{งาน พลังงาน และหลักการเครื่องปั่นเหวี่ยง}
\index{งานและพลังงาน}\index{เครื่องปั่นเหวี่ยง}\index{แรงเหวี่ยงสัมพัทธ์ RCF}\index{สมดุลโรเตอร์}

\section{งาน พลังงานกล และกำลังของเครื่องมือห้องปฏิบัติการ}

\deepskinfig{figures/sec_3d/ch04_s1.jpg}{การเปลี่ยนพลังงานไฟฟ้าเป็นงานกลและพลังงานจลน์การหมุนของโรเตอร์ในเครื่องมือห้องปฏิบัติการจุลชีววิทยา}

ในทางฟิสิกส์ \textbf{งาน} (\textbf{Work}, $W$) เกิดขึ้นเมื่อมีแรง $F$ กระทำต่อวัตถุแล้วทำให้วัตถุเคลื่อนที่ไปตามแนวแรงเป็นระยะทาง $s$ โดยมีหน่วยเป็น\textbf{จูล} ($\text{J}$ หรือ $\text{N}\cdot\text{m}$) ส่วน\textbf{พลังงาน} (\textbf{Energy}) คือความสามารถในการทำงาน ซึ่งในระบบกลศาสตร์ประกอบด้วย \textbf{พลังงานจลน์} ($E_k = \frac{1}{2}mv^2$ พลังงานของวัตถุที่กำลังเคลื่อนที่) และ\textbf{พลังงานศักย์โน้มถ่วง} ($E_p = mgh$ พลังงานที่สะสมตามระดับความสูง $h$) สำหรับเครื่องมือวิทยาศาสตร์ทุกชนิดในห้องแล็บจุลชีววิทยา เราใช้ปริมาณที่เรียกว่า \textbf{กำลัง} (\textbf{Power}, $P$) เพื่อบอกอัตราการใช้พลังงานต่อหนึ่งหน่วยเวลา มีหน่วยเป็น\textbf{วัตต์} ($\text{W}$ หรือ $\text{J/s}$) ดังแสดงใน\mbox{ภาพที่} 4.1 และสมการที่ (4.1)

\begin{equation}
    W = F s \qquad \text{และ} \qquad P = \frac{W}{t} = \frac{E}{t}
\end{equation}

\section{ฟิสิกส์ของเครื่องปั่นเหวี่ยง การแปลงค่า RPM เป็น RCF}

\deepskinfig{figures/sec_3d/ch04_s2.jpg}{โครงสร้างหัวปั่นโรเตอร์มุมคงที่ การจัดเรียงหลอดถ่วงสมดุลอย่างสมมาตร และการอัดแน่นของเม็ดตะกอนเซลล์ (Pellet) ภายใต้แรงเหวี่ยงสัมพัทธ์ RCF ($\times g$)}

ดังที่ได้ศึกษาในบทที่ 3 เซลล์แบคทีเรียมีขนาดเล็กมากจนการปล่อยให้ตกตะกอนด้วยแรงโน้มถ่วงโลกเพียงอย่างเดียว ($1\times g$) ต้องใช้เวลานานหลายวัน ดังนั้น ห้องปฏิบัติการจุลชีววิทยาจึงต้องใช้ \textbf{เครื่องปั่นเหวี่ยง} (\textbf{Centrifuge}) เพื่อสร้างความเร่งสู่ศูนย์กลางที่มีค่าสูงกว่าแรงโน้มถ่วงโลกหลายพันเท่า ปริมาณที่บอกให้ทราบว่าแรงเหวี่ยงที่กระทำต่อเซลล์มีค่าเป็นกี่เท่าของแรงโน้มถ่วงโลก เรียกว่า \textbf{แรงเหวี่ยงสัมพัทธ์} (\textbf{Relative Centrifugal Force} หรือ \textbf{RCF} นิยมเขียนหน่วยเป็น $\times g$) ซึ่งคำนวณได้จากรัศมีของหัวปั่นโรเตอร์ในหน่วยเซนติเมตร ($r_{\text{cm}}$) และความเร็วรอบต่อนาที (\textbf{Revolutions Per Minute}, $\text{RPM}$) ตามสมการมาตรฐานที่ (4.2)

\begin{equation}
    \text{RCF} = 1.118 \times 10^{-5} \times r_{\text{cm}} \times (\text{RPM})^2
\end{equation}

\begin{conceptbox}[ทำไมห้ามบอกแค่ค่า RPM เวลาเขียนรายงานวิจัยจุลชีววิทยา]
จากสมการที่ (4.2) แม้เราจะตั้งความเร็วรอบ $\text{RPM}$ เท่ากัน แต่ถ้าเครื่องปั่นเหวี่ยง 2 เครื่องมี\textbf{รัศมีโรเตอร์ $r_{\text{cm}}$ ไม่เท่ากัน} ค่าแรงเหวี่ยงสัมพัทธ์ ($\text{RCF}$) ที่กระทำต่อเซลล์แบคทีเรียจะไม่เท่ากันเลย ดังนั้น มาตรฐานสากลจึงกำหนดให้ระบุเป็นค่า $\times g$ เสมอ เช่น \textbf{ปั่นเหวี่ยงเก็บเซลล์แบคทีเรียที่ $5{,}000 \times g$ เป็นเวลา 10 นาที}
\end{conceptbox}

นอกจากนี้ ขณะที่โรเตอร์หมุนด้วยความเร็วสูงถึง $10{,}000\text{--}15{,}000\text{ rpm}$ หลอดตัวอย่างแต่ละหลอดจะมีแรงหนีศูนย์กลางขนาดมหาศาลดึงออกไปตามแนวรัศมี ดังนั้น นักศึกษา\textbf{ต้องวางหลอดที่มีมวลเท่ากันให้อยู่ในช่องตรงข้ามกันอย่างสมมาตรเสมอ} เพื่อให้ผลรวมเวกเตอร์ของแรงหนีศูนย์กลางที่จุดหมุนมีค่าเป็นศูนย์ ($\sum \vect{F}_c = 0$) ดังแสดงใน\mbox{ภาพที่} 4.2 และ\mbox{ภาพที่} 4.3 หากวางหลอดไม่สมดุล แกนโรเตอร์จะสั่นอย่างรุนแรงและอาจแตกหักก่อให้เกิดอันตรายร้ายแรงและทำให้ละอองเชื้อโรคฟุ้งกระจาย

\begin{figure}[H]
\centering
\resizebox{0.86\textwidth}{!}{%
\begin{tikzpicture}[font=\sffamily, >=Stealth]
    % Balanced Rotor (Left)
    \begin{scope}[shift={(0,0)}]
        \fill[boxGreenBg] (0,0) circle (2.4cm);
        \draw[bioEmerald, line width=1.8pt] (0,0) circle (2.4cm);
        \fill[charcoal] (0,0) circle (0.35cm);
        \foreach \ang in {0, 60, 120, 180, 240, 300} {
            \draw[physSlate, line width=1.0pt, fill=white] (\ang:1.55cm) circle (0.36cm);
        }
        % 2 Balanced Tubes (0 and 180)
        \fill[bioEmerald] (0:1.55cm) circle (0.32cm);
        \fill[bioEmerald] (180:1.55cm) circle (0.32cm);
        \draw[->, line width=1.5pt, color=physRoyalBlue] (0:0.4cm) -- (0:1.15cm);
        \draw[->, line width=1.5pt, color=physRoyalBlue] (180:0.4cm) -- (180:1.15cm);
        \node[font=\bfseries\color{bioEmerald!90!black}] at (0,-3.0) {(ก) สมดุล 2 หลอดตรงข้าม ($\sum \vect{F} = 0$)};
    \end{scope}

    % Balanced 3 Tubes (Middle)
    \begin{scope}[shift={(6.0,0)}]
        \fill[boxGreenBg] (0,0) circle (2.4cm);
        \draw[bioEmerald, line width=1.8pt] (0,0) circle (2.4cm);
        \fill[charcoal] (0,0) circle (0.35cm);
        \foreach \ang in {0, 60, 120, 180, 240, 300} {
            \draw[physSlate, line width=1.0pt, fill=white] (\ang:1.55cm) circle (0.36cm);
        }
        % 3 Balanced Tubes (0, 120, 240)
        \fill[physCyan] (0:1.55cm) circle (0.32cm);
        \fill[physCyan] (120:1.55cm) circle (0.32cm);
        \fill[physCyan] (240:1.55cm) circle (0.32cm);
        \draw[physCyan, dashed, line width=1.2pt] (0:1.55cm) -- (120:1.55cm) -- (240:1.55cm) -- cycle;
        \node[font=\bfseries\color{bioEmerald!90!black}] at (0,-3.0) {(ข) สมดุล 3 หลอดแบบสามเหลี่ยมด้านเท่า};
    \end{scope}

    % Unbalanced (Right)
    \begin{scope}[shift={(12.0,0)}]
        \fill[boxRedBg] (0,0) circle (2.4cm);
        \draw[crimsonAlert, line width=1.8pt] (0,0) circle (2.4cm);
        \fill[charcoal] (0,0) circle (0.35cm);
        \foreach \ang in {0, 60, 120, 180, 240, 300} {
            \draw[physSlate, line width=1.0pt, fill=white] (\ang:1.55cm) circle (0.36cm);
        }
        \fill[crimsonAlert] (0:1.55cm) circle (0.32cm);
        \fill[crimsonAlert] (60:1.55cm) circle (0.32cm);
        \draw[crimsonAlert, line width=2.2pt] (-1.6,-1.6) -- (1.6,1.6);
        \draw[crimsonAlert, line width=2.2pt] (-1.6,1.6) -- (1.6,-1.6);
        \node[font=\bfseries\color{crimsonAlert}] at (0,-3.0) {(ค) ไม่สมดุล (อันตราย แกนหมุนพัง!)};
    \end{scope}
\end{tikzpicture}%
}
\caption{ผังการจัดเรียงหลอดตัวอย่างในหัวปั่นโรเตอร์เพื่อรักษาสมดุลของแรงหนีศูนย์กลาง}
\end{figure}

\vspace{4pt}
\arucomarkerbadge{3}{4}{เครื่องปั่นเหวี่ยงความเร็วสูง การแปลงค่า RCF (xg) และการจัดสมดุลโรเตอร์}{จำลองการหมุนของโรเตอร์มุมคงที่ การวางหลอดตัวอย่างถ่วงสมดุลตรงข้ามกัน และการอัดแน่นของเม็ดตะกอนเซลล์ (Pellet) ที่ก้นหลอด}{ความเร็วรอบ $\text{RPM}$ ($1{,}000\text{--}14{,}000$), รัศมีโรเตอร์ $r_{\text{cm}}$ ($\text{cm}$), สมดุลหลอด}


\begin{worked}[การคำนวณค่าแรงเหวี่ยงสัมพัทธ์ RCF ของเครื่องปั่นเหวี่ยง]
\textbf{โจทย์ปัญหา}\quad นักศึกษาใช้เครื่องปั่นเหวี่ยงที่มีรัศมีโรเตอร์ $r_{\text{cm}} = 10\text{ cm}$ เพื่อปั่นเก็บตะกอนเซลล์แบคทีเรีย \textit{E. coli} โดยตั้งความเร็วรอบที่ $\text{RPM} = 5{,}000\text{ รอบต่อนาที}$ จงคำนวณหาค่าแรงเหวี่ยงสัมพัทธ์ ($\text{RCF}$) ในหน่วย $\times g$

\textbf{วิธีทำทีละขั้นตอนฉบับเข้าใจง่าย}
\begin{enumerate}
    \item \textbf{ขั้นที่ 1 ยกกำลังสองค่าความเร็วรอบ $(\text{RPM})^2$ ก่อน}
    \[
        (\text{RPM})^2 = (5{,}000)^2 = 5{,}000 \times 5{,}000 = 25{,}000{,}000 = 2.5 \times 10^7
    \]
    \item \textbf{ขั้นที่ 2 แทนค่าลงในสูตร $\text{RCF} = 1.118 \times 10^{-5} \times r_{\text{cm}} \times (\text{RPM})^2$}
    \[
        \text{RCF} = 1.118 \times 10^{-5} \times 10 \times (2.5 \times 10^7)
    \]
    \item \textbf{ขั้นที่ 3 คูณตัวเลขเข้าด้วยกัน ($10^{-5} \times 10 \times 10^7 = 10^3 = 1{,}000$)}
    \[
        \text{RCF} = 1.118 \times 2.5 \times 1{,}000 = 2{,}795 \times g
    \]
\end{enumerate}
\textbf{คำตอบ}\quad แรงเหวี่ยงสัมพัทธ์ที่กระทำต่อเซลล์แบคทีเรียมีค่าเท่ากับ $2{,}795 \times g$ (ประมาณ $2{,}800$ เท่าของแรงโน้มถ่วงโลก)
\end{worked}

\section{ปฏิบัติการประจำสัปดาห์ที่ 4}

\deepskinfig{figures/sec_3d/ch04_s3.jpg}{ชุดปฏิบัติการที่ 4 สำหรับการวัดรัศมีโรเตอร์ ($r_{\text{cm}}$) คำนวณค่า RCF และสังเกตการแยกชั้นระหว่างตะกอนเซลล์ยีสต์ (Pellet) กับส่วนใส (Supernatant)}

\begin{labbox}[ปฏิบัติการที่ 4\quad การอนุรักษ์พลังงานและการคำนวณ RCF เครื่อง Centrifuge (3 ชั่วโมง)]
\noindent\textbf{ผลลัพธ์การเรียนรู้ที่ประเมิน}\quad CLO1 (พลังงานและแรงเหวี่ยง), CLO2 (การตั้งค่าเครื่อง Centrifuge), CLO3 (ความปลอดภัยในการจัดสมดุลโรเตอร์), CLO4 (รายงานกลุ่ม)

\vspace{4pt}
\noindent\textbf{1. วัสดุและอุปกรณ์ (ดังแสดงใน\mbox{ภาพที่} 4.4)}
\begin{itemize}
    \item เครื่องปั่นเหวี่ยงขนาดตั้งโต๊ะ (Tabletop Centrifuge) และเครื่องปั่นเหวี่ยงหลอดขนาดเล็ก (Microcentrifuge)
    \item ไม้บรรทัดสำหรับวัดรัศมีโรเตอร์ ($r_{\text{cm}}$), หลอดไมโครทิวบ์ขนาด $1.5\text{ mL}$, ไมโครปิเปต และสารแขวนลอยยีสต์อบแห้งในน้ำ
\end{itemize}

\noindent\textbf{2. ขั้นตอนการทดลอง}
\begin{enumerate}
    \item ใช้ไม้บรรทัดวัดรัศมีโรเตอร์ ($r_{\text{cm}}$) จากจุดกึ่งกลางแกนหมุนไปถึงก้นช่องใส่หลอดของเครื่องปั่นเหวี่ยงทั้ง 2 เครื่อง
    \item คำนวณค่า $\text{RCF}$ ที่ความเร็วรอบ $2{,}000\text{ rpm}$, $4{,}000\text{ rpm}$ และ $6{,}000\text{ rpm}$ เปรียบเทียบกับค่าที่แสดงบนหน้าจอดิจิทัลของเครื่อง
    \item ดูดสารแขวนลอยยีสต์ปริมาตร $1{,}000\text{ }\mu\text{L}$ ใส่หลอดไมโครทิวบ์ เตรียมหลอดถ่วงดุลน้ำกลั่นให้มีมวลเท่ากัน วางในตำแหน่งตรงข้ามกัน แล้วปั่นที่ $3{,}000\text{ rpm}$ เป็นเวลา $2$ นาที สังเกตเม็ดตะกอนเซลล์ (\textbf{Pellet}) และส่วนใส (\textbf{Supernatant})
\end{enumerate}

\noindent\textbf{3. ตารางบันทึกผลการทดลอง}
\begin{center}
\footnotesize\setlength{\tabcolsep}{3.5pt}
\begin{tabularx}{\linewidth}{c c c c >{\raggedright\arraybackslash}X}
\toprule
\textbf{รอบ (RPM)} & \textbf{รัศมี $r_{\text{cm}}$} & \textbf{RCF คำนวณ ($\times g$)} & \textbf{RCF หน้าจอ} & \textbf{ลักษณะตะกอนเซลล์} \\
\midrule
$2{,}000\text{ rpm}$ & ............ & ............ & ............ & \dotfill \\
$4{,}000\text{ rpm}$ & ............ & ............ & ............ & \dotfill \\
$6{,}000\text{ rpm}$ & ............ & ............ & ............ & \dotfill \\
\bottomrule
\end{tabularx}
\end{center}
\end{labbox}

\section*{คำถามท้ายบทที่ 4 (ทบทวนแนวคิดและการคำนวณเบื้องต้น)}
\begin{enumerate}
    \item เครื่องปั่นเหวี่ยงขนาดเล็กมีรัศมีโรเตอร์ $r_{\text{cm}} = 8\text{ cm}$ หากตั้งความเร็วรอบไว้ที่ $10{,}000\text{ rpm}$ จะเกิดแรงเหวี่ยงสัมพัทธ์ ($\text{RCF}$) กี่เท่าของแรงโน้มถ่วงโลก ($\times g$)
    \item หากนักศึกษามีหลอดตัวอย่างเชื้อแบคทีเรียเพียง 1 หลอดที่ต้องการปั่นเหวี่ยง นักศึกษาจะต้องปฏิบัติอย่างไรก่อนกดปุ่มเดินเครื่อง เพื่อป้องกันอันตรายต่อแกนโรเตอร์
\end{enumerate}
